% Part: history
% Chapter: biographies 
% Section: alan-turing 

\documentclass[../../../include/open-logic-section]{subfiles}

\begin{document}

\olfileid{his}{bio}{tur} 

\olsection{એલન ટ્યુરિંગ}

એલન ટ્યુરિંગનો જન્મ 23 જૂન
1912ના રોજ લંડનના માઇડા
વેલમાં થયો હતો. તેમને
સૈદ્ધાંતિક સંગણકવિજ્ઞાનના
પિતા માનવામાં આવે છે.
ભૌતિક વિજ્ઞાનો અને ગણિતમાં
તેમનો રસ નાનપણથી જ હતો.
પરંતુ બાળપણમાં તેમની
શાળાઓમાં સાહિત્ય અને
પ્રાચીન વિદ્યાને ભાર
આપાતો હોવાથી તેમની
રુચિઓને ત્યાં પૂરતું
સ્થાન મળતું ન હતું.
આથી શાળામાં તેમનો
દેખાવ નબળો રહ્યો અને
ઘણા શિક્ષકો તેમને
ઠપકો આપતા.

\olphoto{turing-alan}{એલન ટ્યુરિંગ}

ટ્યુરિંગે સ્નાતક તરીકે
કેમ્બ્રિજની કિંગ્સ
કૉલેજમાં ગણિત ભણ્યું.
1936માં સંગણનીય
વિધેયની કલ્પનાને
ચોક્કસ રીતે વ્યાખ્યાયિત
કરવા અને નિર્ણય
સમસ્યાની અનિર્ણેયતા
સાબિત કરવા તેમણે
આજે ટ્યુરિંગ મશીન
કહેવાતું મશીન વિકસાવ્યું.
અલોન્ઝો ચર્ચે પોતાના
લૅમ્બડા કલન વડે આ
પરિણામ અગાઉ સાબિત
કરી દીધું હતું.
તેમ છતાં ચર્ચના
પરિણામનો સંદર્ભ
આપીને ટ્યુરિંગનો
લેખ પ્રકાશિત થયો.
ચર્ચે ટ્યુરિંગને
પ્રિન્સ્ટન બોલાવ્યા,
જ્યાં તેઓ 1936--1938
રહ્યા અને ચર્ચના
માર્ગદર્શન હેઠળ
ડૉક્ટરેટ મેળવી.

તર્કશાસ્ત્રમાં રસ હોવા
છતાં ટ્યુરિંગનો ભૌતિક
વિજ્ઞાનોમાં પહેલાનો
રસ યથાવત્ રહ્યો.
બીજા વિશ્વયુદ્ધ દરમિયાન
બ્લેચલી પાર્ક ખાતેના
બ્રિટિશ ગુપ્તલિપિ-વિશ્લેષણ
વિભાગમાં તેમની વ્યવહારુ
કુશળતા કામે લાગી.
જર્મન નૌકાદળના
સંદેશોમાં વપરાતી
એનિગ્મા ગુપ્તલિપિ
ઉકેલવામાં ટ્યુરિંગે
મુખ્ય ભૂમિકા ભજવી.
આંકડાશાસ્ત્ર અને
ગુપ્તલિપિશાસ્ત્રમાં
તેમની નિપુણતા,
સાથે ઇલેક્ટ્રૉનિક
યંત્રોના પ્રવેશથી,
ટીમ “બૉમ્બ”
નામનું ગુપ્તલિપિ
ઉકેલતું મશીન બનાવી
કોડ ઉકેલી શકી.
ટ્યુરિંગના વિચારો
વિશ્વના પ્રથમ
પ્રોગ્રામ કરી શકાય
તેવા ઇલેક્ટ્રૉનિક
કમ્પ્યુટર કોલૉસસની
રચનામાં પણ મદદરૂપ
થયા. બ્લેચલી
પાર્કમાં તેનો
ઉપયોગ જર્મન
લોરેન્ઝ ગુપ્તલિપિ
ઉકેલવા થયો.

ટ્યુરિંગ સમલૈંગિક
હતા. તેમ છતાં
1942માં તેમણે
બ્લેચલી પાર્કની
પોતાની સહકર્મી
જોન ક્લાર્કને
લગ્નનો પ્રસ્તાવ
મૂક્યો. પાછળથી
તેમણે સગાઈ તોડી
અને પોતે સમલૈંગિક
છે એવું તેમને
કહ્યું. જીવનમાં
તેમના અનેક
પ્રેમસંબંધો રહ્યા,
જોકે તે સમયે
બ્રિટનમાં સમલૈંગિક
કૃત્યો ગુના
ગણાતા હતા.
1952માં તેમના
તે સમયના સાથીના
એક મિત્રે તેમના
ઘરમાં ચોરી કરી.
પોલીસ ફરિયાદ
નોંધાવતાં ટ્યુરિંગે
સમલૈંગિક સંબંધ
હોવાની કબૂલાત
કરી; તેમને
લાગતું હતું કે
સરકાર આવાં
કૃત્યોને કાયદેસર
કરવાની તૈયારીમાં
છે. હકીકતમાં એવું
ન હતું, અને
તેમના પર
“ઘોર અશિષ્ટતા”નો
આરોપ મૂકાયો.
જેલમાં જવાને
બદલે ટ્યુરિંગે
જાતીય ઇચ્છા
ઘટાડતી હોર્મોન
સારવાર પસંદ કરી.
8 જૂન 1954ના
રોજ તેઓ
સાઇનાઇડની
અતિમાત્રાને કારણે
મૃત હાલતમાં
મળ્યા—સંભવતઃ
તે આત્મહત્યા
હતી. 2013માં
મહારાણી એલિઝાબેથ
દ્વિતીયે તેમને
રાજક્ષમા આપી.

\begin{reading}
એલન ટ્યુરિંગના
વિસ્તૃત જીવનચરિત્ર
માટે \citet{Hodges2014}
જુઓ. તેમના જીવન
અને કાર્યથી
પ્રેરાયેલું
\emph{Breaking the Code}
નાટક 1996માં
દૂરદર્શન માટે
બનાવાયું હતું;
તેમાં ડેરેક
જેકોબીએ ટ્યુરિંગની
ભૂમિકા ભજવી.
ઍકૅડેમી પુરસ્કાર
માટે નામાંકિત
\emph{The Imitation Game}
ચલચિત્રમાં બેનેડિક્ટ
કમ્બરબૅચ અને
કીરા નાઇટ્લી
છે; તે પણ
ટ્યુરિંગના જીવન
અને બ્લેચલી
પાર્કમાં તેમના
સમય પર છૂટથી
આધારિત છે
\citep{Imitation2014}.

\citet{Radiolab2012}માં
ટ્યુરિંગના જીવન
અને કાર્ય વિશે
ઘણા પૉડકાસ્ટ
છે. બીબીસીના
હોરાઇઝન કાર્યક્રમનું
\emph{The Strange Life and Death of
  Dr.~Turing}
દસ્તાવેજી ચલચિત્ર
ઑનલાઇન જોઈ શકાય
છે \citep{Sykes1992}.
\citep{Theelen2012}માં
કામ કરતા લેગો
ટ્યુરિંગ મશીનનો
ટૂંકો વિડિયો છે;
2012માં ટ્યુરિંગની
જન્મશતાબ્દી નિમિત્તે
તે બનાવાયું હતું.

ટ્યુરિંગ મશીનો
અને નિર્ણય સમસ્યા
વિશેનો તેમનો મૂળ
લેખ \citet{Turing1937}
છે.
\end{reading}

\end{document}
